\documentclass[twocolumn]{article}
\usepackage[margin=1in]{geometry}
\usepackage{graphicx}
\usepackage{absfigure}

% Filler paragraph command (lipsum replacement)
\newcommand{\fillerpar}{%
  Lorem ipsum dolor sit amet, consectetur adipiscing elit. Sed do
  eiusmod tempor incididunt ut labore et dolore magna aliqua. Ut enim
  ad minim veniam, quis nostrud exercitation ullamco laboris nisi ut
  aliquip ex ea commodo consequat. Duis aute irure dolor in
  reprehenderit in voluptate velit esse cillum dolore eu fugiat nulla
  pariatur. Excepteur sint occaecat cupidatat non proident, sunt in
  culpa qui officia deserunt mollit anim id est laborum.

}

\begin{document}

\title{Test Document for \textsf{absfigure}}
\author{Test}
\maketitle
\listoffigures
\listoftables

%% ============================================================
%% Define absolute figures up front (before their target pages)
%% ============================================================

% Figure on page 2, top of column 1
\begin{absfigure}{page=2, pos=t, col=1}
  \centering
  \rule{0.9\columnwidth}{3cm}% placeholder for \includegraphics
  \caption{Absolute figure: page~2, top, column~1}
  \label{fig:p2t1}
\end{absfigure}

% Figure on page 2, bottom of column 2 (using left/right/top/bottom aliases)
\begin{absfigure}{page=2, pos=bottom, col=right}
  \centering
  \rule{0.9\columnwidth}{2.5cm}
  \caption{Absolute figure: page~2, bottom, right}
  \label{fig:p2b2}
\end{absfigure}

% Spanning figure on page 3, top
\begin{absfigure*}{page=3, pos=t}
  \centering
  \rule{0.9\textwidth}{4cm}
  \caption{Spanning absolute figure: page~3, top}
  \label{fig:p3span}
\end{absfigure*}

%% ============================================================
%% Stacking test: two figures targeting the same page/col/pos.
%% The first-defined figure should have precedence (appear closer
%% to the bottom edge, matching standard figure[b] stacking).
%% ============================================================

% First bottom figure on page 4, col 2 — should be closest to text
\begin{absfigure}{page=4, pos=b, col=2}
  \centering
  \rule{0.9\columnwidth}{2cm}
  \caption{Bottom stack A (defined first) --- page~4, col~2}
  \label{fig:stackA}
\end{absfigure}

% Second bottom figure on page 4, col 2 — should stack below A
\begin{absfigure}{page=4, pos=b, col=2}
  \centering
  \rule{0.9\columnwidth}{1.5cm}
  \caption{Bottom stack B (defined second) --- page~4, col~2}
  \label{fig:stackB}
\end{absfigure}

% Two top figures on the same page/col to test top stacking too
\begin{absfigure}{page=4, pos=t, col=1}
  \centering
  \rule{0.9\columnwidth}{2cm}
  \caption{Top stack A (defined first) --- page~4, col~1}
  \label{fig:topA}
\end{absfigure}

\begin{absfigure}{page=4, pos=t, col=1}
  \centering
  \rule{0.9\columnwidth}{1.5cm}
  \caption{Top stack B (defined second) --- page~4, col~1}
  \label{fig:topB}
\end{absfigure}

% Spanning figure at bottom of page 3
\begin{absfigure*}{page=3, pos=bottom}
  \centering
  \rule{0.9\textwidth}{3cm}
  \caption{Spanning absolute figure: page~3, bottom}
  \label{fig:p3spanbot}
\end{absfigure*}

% Figure placed on the last page via page=end
\begin{absfigure}{page=end}
  \centering
  \rule{0.9\columnwidth}{3cm}
  \caption{Figure placed at the end of the document}
  \label{fig:end}
\end{absfigure}

%% ============================================================
%% Table tests (abstable and abstable*)
%% ============================================================

% Table on page 2, bottom of column 1
\begin{abstable}{page=2, pos=b, col=1}
  \centering
  \caption{Absolute table: page~2, bottom, column~1}
  \label{tab:p2b1}
  \begin{tabular}{lcc}
    \hline
    Item & Value & Unit \\
    \hline
    Alpha & 1.0 & m/s \\
    Beta  & 2.5 & m/s \\
    \hline
  \end{tabular}
\end{abstable}

% Spanning table on page 3, top
\begin{abstable*}{page=3, pos=t}
  \centering
  \caption{Spanning absolute table: page~3, top}
  \label{tab:p3span}
  \begin{tabular}{lccccc}
    \hline
    Experiment & Trial~1 & Trial~2 & Trial~3 & Mean & Std \\
    \hline
    Control    & 4.2     & 4.5     & 4.1     & 4.27 & 0.21 \\
    Treatment  & 6.1     & 5.8     & 6.3     & 6.07 & 0.25 \\
    \hline
  \end{tabular}
\end{abstable*}

% Table placed at the end via page=end
\begin{abstable}{page=end}
  \centering
  \caption{Table placed at the end of the document}
  \label{tab:end}
  \begin{tabular}{lr}
    \hline
    Category & Count \\
    \hline
    Placed   & 100\% \\
    \hline
  \end{tabular}
\end{abstable}

%% ============================================================
%% page=current test: places on whatever page is being built
%% when this code is processed (should be page 1)
%% ============================================================

\begin{absfigure}{page=current, pos=b, col=2}
  \centering
  \rule{0.9\columnwidth}{2cm}
  \caption{Figure via page=current (should be page~1)}
  \label{fig:current}
\end{absfigure}

% Default page (omitted = current page, should also be page 1)
\begin{absfigure}{pos=t, col=2}
  \centering
  \rule{0.9\columnwidth}{1.5cm}
  \caption{Figure with default page (should be page~1)}
  \label{fig:default}
\end{absfigure}

%% ============================================================
%% Normal figure (should still work alongside absolute ones)
%% ============================================================

\begin{figure}[t]
  \centering
  \rule{0.9\columnwidth}{2cm}
  \caption{A normal LaTeX float figure}
  \label{fig:normal}
\end{figure}

%% ============================================================
%% Cross-references (will resolve after 2 compilations)
%% ============================================================

\section{Introduction}

This document tests the \textsf{absfigure} package.
See Figure~\ref{fig:p2t1} on page~\pageref{fig:p2t1},
Figure~\ref{fig:p2b2} on page~\pageref{fig:p2b2},
Figure~\ref{fig:p3span} on page~\pageref{fig:p3span},
Figure~\ref{fig:p3spanbot} on page~\pageref{fig:p3spanbot},
Figure~\ref{fig:end} on page~\pageref{fig:end},
Figure~\ref{fig:current} on page~\pageref{fig:current},
and Figure~\ref{fig:default} on page~\pageref{fig:default}.
The normal float is Figure~\ref{fig:normal}.

Table tests:
Table~\ref{tab:p2b1} on page~\pageref{tab:p2b1},
Table~\ref{tab:p3span} on page~\pageref{tab:p3span},
and Table~\ref{tab:end} on page~\pageref{tab:end}.

Stacking test (page~4): bottom col~2 has
Figures~\ref{fig:stackA} and~\ref{fig:stackB};
top col~1 has Figures~\ref{fig:topA} and~\ref{fig:topB}.

%% ============================================================
%% Filler text to generate multiple pages
%% ============================================================

\fillerpar \fillerpar \fillerpar \fillerpar
\fillerpar \fillerpar \fillerpar \fillerpar

\section{More Content}

\fillerpar \fillerpar \fillerpar \fillerpar
\fillerpar \fillerpar \fillerpar \fillerpar

\section{Even More Content}

\fillerpar \fillerpar \fillerpar \fillerpar
\fillerpar \fillerpar \fillerpar \fillerpar

\section{Final Section}

\fillerpar \fillerpar \fillerpar \fillerpar
\fillerpar \fillerpar \fillerpar \fillerpar

\end{document}
